\documentclass[12pt]{article}
\usepackage{seantex}

% --- Document Info ---
\title{Linear Algebra II: Serie 22}
\author{Sean Findlay}
\date{\today}

% --- Start Document ---
\begin{document}

\maketitle

\begin{problem}{}{}
    For which values of $a \in \R$ does the expression
    $$
        \langle x, y \rangle := x_1y_1 + ax_1y_2 + ax_2y_1 + 7x_2y_2
    $$
    where $x = \begin{pmatrix} x_1 \\ x_2 \end{pmatrix}$ and $y = \begin{pmatrix} y_1 \\ y_2 \end{pmatrix}$, define an inner product on $\R^2$?
\end{problem}

We must check for which $a \in \R$ the following properties hold:
\begin{enumerate}
    \item \textbf{Linearity in the first variable}: Let $x = \begin{pmatrix} x_1 \\ x_2 \end{pmatrix}$, $x' = \begin{pmatrix} x_1' \\ x_2' \end{pmatrix}$, $y = \begin{pmatrix} y_1 \\ y_2 \end{pmatrix} \in \R^2$. It follows that
    $$
        \langle x + x', y \rangle = (x_1 + x_1')y_1 + a(x_1 + x_1')y_2 + a(x_2 + x_2')y_1 + 7(x_2 + x_2')y_2
    $$
    $$
        = x_1y_1 +ax_1y_2 + ax_2y_1 + 7x_2y_2 + x_1'y_1 +ax_1'y_2 + ax_2'y_1 + 7x_2'y_2
    $$
    $$
        = \langle x, y \rangle + \langle x', y \rangle
    $$

    Now, let $\alpha \in \R$. We have that
    $$
        \langle \alpha x, y \rangle = \alpha x_1y_1 + a\alpha x_1y_2 + a\alpha x_2y_1 + 7\alpha x_2y_2
    $$
    $$
        = \alpha (x_1y_1 + ax_1y_2 + ax_2y_1 + 7x_2y_2) = \alpha \langle x, y \rangle
    $$

    So we have linearity in the first variable regardless of the choice of $a \in \R$.

    \item \textbf{Symmetry}: Swapping the roles of $x$ and $y$ in the definition of $\langle x, y \rangle$ gives us precisely $\langle y, x \rangle$. So symmetry is guaranteed regardless of the choice of $a \in \R$.
    \item \textbf{Linearity in the second variable}: We can show linearity in the second variable by first applying symmetry, then linearity in the first variable, then swapping back again using symmetry.
    \item \textbf{Positivity}: Let $x = \begin{pmatrix} x_1 \\ x_2 \end{pmatrix}$, then we have
    $$
        \langle x, x \rangle = x_1x_1 + ax_1x_2 + ax_2x_1 + 7x_2x_2 = x_1^2 + 2ax_1x_2 + 7x_2^2
    $$
    $$
        = (x_1^2 + 2ax_1x_2 + a^2x_2^2) - a^2x_2^2 + 7x_2^2 = (x_1 + ax_2)^2 + (7-a^2) x_2^2
    $$

    If $x \neq 0$, then
    \begin{itemize}
        \item either $x_2 = 0$: then $x_1$ must be non-zero for the expression to be strictly positive, which is given because $x \neq 0$.
        \item or $x_2 \neq 0$: then the first term $(x_1 + ax_2)^2$ is $\geq 0$. For the entire term to then be strictly positive, we must require that $(7-a^2)x_2^2 > 0$:
            $$
                7 - a^2 > 0 \iff 7 > a^2 \iff \pm \sqrt{7} > a \iff a \in (-\sqrt{7}, \sqrt{7})
            $$
    \end{itemize}
\end{enumerate}

So the expression defines an inner product iff $a \in (-\sqrt{7}, \sqrt{7})$.

\begin{problem}{}{}
    Let $V$ be the vector space of real polynomials of degree at most $n$.
    \begin{enumerate}[label=\alph*)]
        \item Show that the expression
        $$
            \langle p, q \rangle := \int_{0}^{\infty} p(t)q(t)e^{-t}\dd t
        $$
        defines an inner product on $V$.
        \item Find the matrix of the inner product with respect to the basis $1, x, \dots, x^n$.
    \end{enumerate}
\end{problem}

\begin{enumerate}[label=\alph*)]
    \item We need to check the four properties of inner products:
    \begin{enumerate}[label=\arabic*.]
        \item \textbf{Linearity in the first variable}: Let $p_1, p_2, q \in \R[x]_{\leq n}$. We have
        $$
            \langle p_1 + p_2, q \rangle = \int_{0}^{\infty} (p_1 + p_2)(t)q(t)e^{-t}\dd t = \int_{0}^{\infty} (p_1(t) + p_2(t))q(t)e^{-t}\dd t
        $$
        $$
            = \int_{0}^{\infty} p_1(t)q(t)e^{-t} + p_2(t)q(t)e^{-t}\dd t = \int_{0}^{\infty} p_1(t)q(t)e^{-t} \dd t + \int_{0}^{\infty} p_2(t)q(t)e^{-t}\dd t
        $$
        $$
            = \langle p_1, q \rangle + \langle p_2, q \rangle
        $$

        Let $\alpha \in \R,\ p, q \in \R[x]_{\leq n}$. Then
        $$
            \langle \alpha p, q \rangle = \int_{0}^{\infty} (\alpha p)(t)q(t)e^{-t}\dd t = \int_{0}^{\infty} \alpha p(t)q(t)e^{-t}\dd t = \alpha \int_{0}^{\infty} p(t)q(t)e^{-t}\dd t
        $$

        $$
            = \alpha \langle p, q \rangle
        $$

        So we have linearity in the first variable.

        \item \textbf{Symmetry}: Let $p, q \in \R[x]_{\leq n}$.
        $$
            \langle q, p \rangle = \int_{0}^{\infty} q(t)p(t)e^{-t}\dd t = \int_{0}^{\infty} p(t)q(t)e^{-t}\dd t = \langle p, q \rangle
        $$

        So symmetry holds.

        \item \textbf{Linearity in the second variable}: We can make use of linearity in the first variable and symmetry. Let $p, q_1, q_2 \in \R[x]_{\leq n}$. We have
        $$
            \langle p, q_1 + q_2 \rangle = \langle q_1 + q_2, p \rangle = \langle q_1, p \rangle + \langle q_2, p \rangle = \langle p, q_1 \rangle + \langle p, q_2 \rangle
        $$

        Let $\alpha \in \R,\ p, q \in \R[x]_{\leq n}$. Then
        $$
            \langle p, \alpha q \rangle = \langle \alpha q, p \rangle = \alpha \langle q, p \rangle = \alpha \langle p, q \rangle
        $$

        So we have linearity in the second variable.

        \item \textbf{Positivity}: Let $p \in \R[x]_{\leq n}$. We have
        $$
            \langle p, p \rangle = \int_{0}^{\infty} p(t)p(t)e^{-t}\dd t = \int_{0}^{\infty} p(t)^2e^{-t}\dd t
        $$

        Note that $p(t)^2 \geq 0 \text{ and } e^{-t} > 0\ \forall t \in [0, \infty)$. Assume $p \not\equiv 0$. Because $p(t)$ is a non-zero polynomial of degree less than or equal to $n$, we know that $p(t)$ has at most $n$ roots. Therefore, $p(t)^2e^{-t}$ is a continuous, non-negative function that is strictly positive except for at finitely many points (its roots). The Riemann integral of such a function over $[0, \infty)$ must be strictly positive. Therefore, $\langle p, p \rangle > 0$.
        
    \end{enumerate}

    \item Let $A \in M_{(n+1) \times (n+1)}(\R)$ be the matrix of this inner product that we are looking for. The entry at $A_{ij}$ is our inner product applied to the $i$th and $j$th basis vector:
    $$
        A_{ij} := \langle x^i, x^j \rangle = \int_{0}^{\infty} t^i t^j e^{-t}\dd t = \int_{0}^{\infty} t^{i + j} e^{-t}\dd t
    $$

    This looks very similar to the definition of the gamma function:
    $$
        \Gamma(z) = \int_{0}^{\infty} t^{z - 1} e^{-t}\dd t \quad \text{where } z \in \C,\ \mathrm{Re}(z) > 0
    $$

    Bring our inner product expression into this form:
    $$
        A_{ij} := \int_{0}^{\infty} t^{i + j} e^{-t}\dd t = \int_{0}^{\infty} t^{(i + j + 1) - 1} e^{-t}\dd t = \Gamma(i + j + 1)
    $$

    $i + j + 1$ is just a real number greater or equal to $1$, so this is just the standard factorial function shifted up by one:
    $$
        A_{ij} := \Gamma(i + j + 1) = (i + j)!
    $$

    So the matrix of this inner product is:
    $$
        A = \begin{pmatrix}
            0! & \dots & n! \\
            \vdots & \ddots & \vdots \\
            n! & \dots & (2n)! \\
        \end{pmatrix}
    $$
\end{enumerate}

\begin{problem}{}{}
    Let $V = \R^2$ endowed with the standard scalar product, and for $i = 1, 2$, let $v_i \in V \setminus \{0\}$. Show that the formula
    $$
        \langle v_1, v_2 \rangle = \lVert v_1 \rVert \lVert v_2 \rVert \cos(\widehat{v_1, v_2})
    $$
    
    defining the cosine of an angle is rotation-invariant. In other words, show that for any rotation of the plane $R: V \rightarrow V$, we have
    $$
        \cos(\widehat{v_1, v_2}) = \cos(R \widehat{v_1, v_2})
    $$
\end{problem}

Let $R$ be an arbitrary rotation of the plane $R: V \rightarrow V$. For rotations, we have that 
$$
    \langle R v_1, R v_2 \rangle = \langle v_1, v_2 \rangle
$$

We can write
$$
    \langle R v_1, R v_2 \rangle = \lVert R v_1 \rVert \lVert R v_2 \rVert \cos(\widehat{R v_1, R v_2}) \iff \cos(\widehat{R v_1, R v_2}) = \frac{\langle R v_1, R v_2 \rangle}{\lVert R v_1 \rVert \lVert R v_2 \rVert}
$$

Note that $\lVert R v_1 \rVert \lVert R v_2 \rVert \neq 0$ by positivity of the norm and the fact that $v_1, v_2 \in V \setminus \{0\}$.

We want to show
$$
    \cos(\widehat{v_1, v_2}) \overset{?}= \cos(R \widehat{v_1, v_2}) \iff \frac{\langle v_1, v_2 \rangle}{\lVert v_1 \rVert \lVert v_2 \rVert} \overset{?}= \frac{\langle R v_1, R v_2 \rangle}{\lVert R v_1 \rVert \lVert R v_2 \rVert}
$$

We already know that $\langle v_1, v_2 \rangle = \lVert v_1 \rVert \lVert v_2 \rVert \cos(\widehat{v_1, v_2})$, so
$$
    \iff \frac{\langle v_1, v_2 \rangle}{\lVert v_1 \rVert \lVert v_2 \rVert} \overset{?}= \frac{\langle v_1, v_2 \rangle}{\lVert R v_1 \rVert \lVert R v_2 \rVert}
$$

We have already shown that the numerators are equal. It remains to show that the denominators are equal:
$$
    \iff \lVert R v_1 \rVert \lVert R v_2 \rVert \overset{?}= \lVert v_1 \rVert \lVert v_2 \rVert \iff \sqrt{\langle R v_1, R v_1 \rangle} \cdot \sqrt{\langle R v_2, R v_2 \rangle} \overset{?}= \sqrt{\langle v_1, v_1 \rangle} \cdot \sqrt{\langle v_2, v_2 \rangle}
$$

As stated earlier, rotations do not change the inner product ($\langle R v_1, R v_2 \rangle = \langle v_1, v_2 \rangle$). So our problem is equivalent to
$$
    \iff \sqrt{\langle v_1, v_1 \rangle} \cdot \sqrt{\langle v_2, v_2 \rangle} \overset{?}= \sqrt{\langle v_1, v_1 \rangle} \cdot \sqrt{\langle v_2, v_2 \rangle}
$$

which is clearly true. So by going through the chain of implications backwards, we can show that, indeed, $\cos(\widehat{v_1, v_2}) = \cos(R \widehat{v_1, v_2})$. \qed

\begin{problem}{}{}
    Let $A \in M_{n \times n}(\R)$. Show that
    \begin{enumerate}[label=\alph*)]
        \item The matrix $A^T A$ is symmetric.
        \item The matrix $A^T A$ is positive definite iff $A$ is invertible.
        \item It holds that $\mathrm{rank}(A^T A) = \mathrm{rank}(A)$.
    \end{enumerate}
\end{problem}

\begin{enumerate}[label=\alph*)]
    \item We want to show that $(A^T A)^T = A^T A$.
    $$
        (A^T A)^T = A^T(A^T)^T = A^T A
    $$

    The transpose of a product of matrices is the product in reverse order with each factor transposed. So we are done. \qed
    \item 
        \begin{itemize}
            \item \underline{$\implies$}: Assume that $A^T A$ is positive definite. Show that $A$ is invertible.
            $A^T A $ is positive definite, so 
            $$
                v^T A^T A v > 0 \quad \forall v \in \R^n,\ v \neq 0 \iff (A v)^T A v > 0 \quad \forall v \in \R^n,\ v \neq 0
            $$
            $$
                \lVert A v \rVert^2 > 0 \quad \forall v \in \R^n,\ v \neq 0
            $$

            Assume $A$ is not invertible. Then, the kernel of $A$ is not trivial, i.e. there exists a non-zero vector $W \in \R^n$ such that $Aw = 0$. But then, $\lVert Aw \rVert^2 = 0$. This contradicts our assumption that $A^T A$ is positive definite, because then we need that $\lVert A v \rVert^2 > 0 \quad \forall v \in \R,\ v \neq 0$.

            Therefore, it must be that $A$ is invertible.
            
            \item \underline{$\impliedby$}: Assume $A$ invertible. We already know from (a) that $A^T A$ is symmetric. Show that $A^T A$ is positive definite.
            $$
                A \text{ invertible } \iff \ker A = \{0\} \iff \lVert A v \rVert^2 > 0 \quad \forall v \in \R^n,\ v \neq 0
            $$
            $$
                \iff (A v)^T A v > 0 \quad \forall v \in \R^n,\ v \neq 0 \iff v^T A^T A v > 0 \quad \forall v \in \R^n,\ v \neq 0 
            $$
            $\iff A^T A$ is positive definite.
        \end{itemize}

    \item Show that $\mathrm{rank}(A^T A) = \mathrm{rank}(A)$.
    $$
        \iff \dim\img(A^T A) = \dim\img(A) \iff \dim\R^n - \dim\ker(A^T A) = \dim\R^n - \dim\ker(A)
    $$
    $$
        \iff \dim\ker(A^T A) = \dim\ker(A)
    $$

    We can show that actually, $\ker(A^T A) = \ker(A)$.

    \begin{itemize}
        \item \underline{$\ker(A) \subseteq \ker(A^T A)$}. Take $x \in \ker(A)$. Then
        $$
            A x = 0 \implies A^T A x = 0 \iff x \in \ker(A^T A)
        $$
        So $x \in \ker(A) \implies x \in \ker(A^T A)$. Therefore, $\ker(A) \subseteq \ker(A^T A)$ holds.
        \item \underline{$\ker(A^T A) \subseteq \ker(A)$}: Take $x \in \ker(A^T A)$. Then
        $$
            A^T A x = 0 \iff x^T A^T A x = x^T \cdot 0 \iff (A x)^T A x = 0
        $$
        $$
            \iff \lVert A x \rVert^2 = 0 \iff \lVert A x \rVert = 0 \iff \sqrt{\langle A x, A x \rangle} = 0
        $$

        By the positivity of the inner product, this implies that $A x$ must be equal to zero. So $x \in \ker(A)$. We have shown that $\ker(A^T A) \subseteq \ker(A)$.
    \end{itemize}

    Because $\ker(A) \subseteq \ker(A^T A)$ and $\ker(A^T A) \subseteq \ker(A)$ both hold, we have shown that $\ker(A^T A) = \ker(A)$. This clearly implies that $\dim\ker(A^T A) = \dim\ker(A)$. By the above chain of implications, this means that $\mathrm{rank}(A^T A) = \mathrm{rank}(A)$. \qed
\end{enumerate}

\begin{problem}{}{}
    \begin{enumerate}[label=\alph*)]
        \item Let $\lVert \cdot \rVert$ be a norm on the $\R$-vector space $V$. Show that the norm is induced by an inner product $\langle \cdot, \cdot \rangle$ on $V$ iff it satisfies the parallelogram identity
        $$
            \lVert x + y \rVert^2 + \lVert x - y \rVert^2 = 2 \lVert x \rVert^2 + 2 \lVert y \rVert^2 \quad \forall x, y \in V
        $$
        \item Let $V$ be a finite-dimensional $\R$ vector space. Consider the following map:
        $$
            \lVert \cdot \rVert_1: V \rightarrow \R_{\geq 0}, \quad v = (v_1,\dots,v_n) \mapsto \sum_{i = 1}^{n} |v_i|
        $$

        Check that $\lVert \cdot \rVert_1$ defines a norm on $V$ and prove that it does not come from a scalar product.
    \end{enumerate}
\end{problem}

\begin{enumerate}[label=\alph*)]
    \item
    \begin{itemize}
        \item \underline{$\implies$}: Assume $\lVert \cdot \rVert$ is induced by an inner product $\langle \cdot, \cdot \rangle$ on $V$. This means
        $$
            \lVert v \rVert = \sqrt{\langle v, v \rangle} \quad \forall v \in V
        $$

        Fix $x, y \in V$. Show that the parallelogram identity holds:
        $$
            \lVert x + y \rVert^2 + \lVert x - y \rVert^2 \overset{?}= 2 \lVert x \rVert^2 + 2 \lVert y \rVert^2
        $$
        $$
            \iff \langle x + y, x + y \rangle + \langle x-y, x-y \rangle \overset{?}= 2 \langle x, x \rangle + 2 \langle y, y \rangle
        $$

        For the LHS, using the bilinearity and symmetry of inner products, we get
        $$
            \langle x + y, x + y \rangle + \langle x-y, x-y \rangle  
        $$
        $$
            = \langle x, x \rangle + \langle x, y \rangle + \langle y, x \rangle + \langle y, y \rangle + \langle x, x \rangle - \langle x, y \rangle - \langle y, x \rangle + \langle y, y \rangle
        $$
        $$
            = 2 \langle x, x \rangle + 2 \langle y, y \rangle
        $$

        Substituting this into the LHS, we get
        $$
            \langle x + y, x + y \rangle + \langle x-y, x-y \rangle \overset{?}= 2 \langle x, x \rangle + 2 \langle y, y \rangle \iff 2 \langle x, x \rangle + 2 \langle y, y \rangle = 2 \langle x, x \rangle + 2 \langle y, y \rangle
        $$

        This holds, so following the chain of implications backwards, we have shown that the parallelogram identity holds.

        \item \underline{$\impliedby$}: Assume the parallelogram identity holds for $\lVert \cdot \rVert$. Show that $\lVert \cdot \rVert$ is induced by an inner product $\langle \cdot, \cdot \rangle$.
        
        Fix $x, y \in V$. By assumption,
        $$
            \lVert x + y \rVert^2 + \lVert x - y \rVert^2 = 2 \lVert x \rVert^2 + 2 \lVert y \rVert^2 \iff \frac{1}{2}\left(\lVert x + y \rVert^2 + \lVert x - y \rVert^2\right) = \lVert x \rVert^2 + \lVert y \rVert^2
        $$

        Using a hint given in the lecture, define
        $$
            \langle u, v \rangle := \frac{1}{2}\left(\lVert u + v \rVert^2 - \lVert u \rVert^2 - \lVert v \rVert^2\right)
        $$

        Show that this definition of $\langle \cdot, \cdot \rangle$ satisfies the properties of an inner product.
        \begin{itemize}
            \item \textbf{Symmetry}: We have
            $$
                \langle v, u \rangle = \frac{1}{2}\left(\lVert v + u \rVert^2 - \lVert v \rVert^2 - \lVert u \rVert^2\right) = \frac{1}{2}\left(\lVert u + v \rVert^2 - \lVert u \rVert^2 - \lVert v \rVert^2\right) = \langle u, v \rangle
            $$
            So symmetry holds.
            \item \textbf{Positivity}: Let $v \in V$ with $v \neq 0$. Show that $\langle v, v \rangle > 0$.
            $$
                \langle v, v \rangle = \frac{1}{2}\left(\lVert v + v \rVert^2 - \lVert v \rVert^2 - \lVert v \rVert^2\right) = \frac{1}{2}\left(\lVert 2 v \rVert^2 - 2 \lVert v \rVert^2\right)
            $$
            $$
                \xRightarrow{homogeneity} \frac{1}{2}\left(4 \lVert v \rVert^2 - 2 \lVert v \rVert^2\right) = 2 \frac{1}{2} \lVert v \rVert^2 = \lVert v \rVert^2
            $$

            By the positivity of the norm, 
            $$
                \lVert v \rVert^2 > 0 \iff \langle v, v \rangle > 0
            $$

            So we have shown positivity of this inner product.
            \item \textbf{Linearity in the first variable}: First, show additivity. Fix $u_1, u_2, v \in V$. We want to show that
            $$
                \langle u_1 + u_2, v \rangle = \langle u_1, v \rangle + \langle u_2, v \rangle
            $$
            Using the definition of our inner product, we have:
            $$
            \begin{aligned}
                \langle u_1, v \rangle + \langle u_2, v \rangle &= \frac{1}{2} \left( \lVert u_1 + v \rVert^2 - \lVert u_1 \rVert^2 - \lVert v \rVert^2 \right)
                + \frac{1}{2} \left( \lVert u_2 + v \rVert^2 - \lVert u_2 \rVert^2 - \lVert v \rVert^2 \right) \\
                &= \frac{1}{2} \left[ \left( \lVert u_1 + v \rVert^2 + \lVert u_2 + v \rVert^2 \right) - \left( \lVert u_1 \rVert^2 + \lVert u_2 \rVert^2 \right) - 2\lVert v \rVert^2 \right]
            \end{aligned}
            $$

            We can now substitute the parallelogram equality in two places:
            $$
            \begin{aligned}
                &= \frac{1}{2} \bigg[ \frac{1}{2} \left( \lVert u_1 + u_2 + 2v \rVert^2 + \lVert u_1 - u_2 \rVert^2 \right)
                - \frac{1}{2} \left( \lVert u_1 + u_2 \rVert^2 + \lVert u_1 - u_2 \rVert^2 \right) - 2\lVert v \rVert^2 \bigg] \\
                &= \frac{1}{4} \left( \lVert u_1 + u_2 + 2v \rVert^2 + \lVert u_1 - u_2 \rVert^2 - \lVert u_1 + u_2 \rVert^2 - \lVert u_1 - u_2 \rVert^2 - 4\lVert v \rVert^2 \right) \\
                &= \frac{1}{4} \left( \lVert u_1 + u_2 + 2v \rVert^2 - \lVert u_1 + u_2 \rVert^2 - 4\lVert v \rVert^2 \right)
            \end{aligned}
            $$

            Use the parallelogram identity again to get the following substitution:
            $$
            \begin{aligned}
                \lVert (u_1 + u_2 + v) + v \rVert^2 + \lVert(u_1 + u_2 + v) - v \rVert^2 = 2 \lVert u_1 + u_2 + v \rVert^2 + 2 \lVert v \rVert^2 \\
                \iff \lVert u_1 + u_2 + 2v \rVert^2 = 2 \lVert u_1 + u_2 + v \rVert^2 + 2 \lVert v \rVert^2 - \lVert u_1 + u_2 \rVert^2 \\
            \end{aligned}
            $$

            Substitute this back into the original expression:
            $$
            \begin{aligned}
                &= \frac{1}{4} \left(2 \lVert u_1 + u_2 + v \rVert^2 + 2 \lVert v \rVert^2 - \lVert u_1 + u_2 \rVert^2 - \lVert u_1 + u_2 \rVert^2 - 4\lVert v \rVert^2 \right) \\
                &= \frac{1}{4} \left(2 \lVert u_1 + u_2 + v \rVert^2 - 2 \lVert v \rVert^2 - 2 \lVert u_1 + u_2 \rVert^2 \right) \\
                &= \frac{1}{2} \left(\lVert (u_1 + u_2) + v \rVert^2 - \lVert u_1 + u_2 \rVert^2 - \lVert v \rVert^2 \right) \\
                &= \langle u_1 + u_2, v \rangle
            \end{aligned}
            $$

            So, we have shown that $\langle u_1, v \rangle + \langle u_2, v \rangle = \langle u_1 + u_2, v \rangle$

            Next, show homogeneity in the first variable. Let $\alpha \in \R,\ u, v \in V$. We want to show
            $$
                \langle \alpha u, v \rangle = \alpha \langle u, v \rangle
            $$

            We can use additivity, which we have already shown.
            $$
                \langle 2 u, v \rangle = \langle u + u, v \rangle = \langle u, v \rangle + \langle u, v \rangle = 2 \langle u, v \rangle
            $$

            We can show by induction that this holds for all natural numbers. Let the above case ($n = 2$) be our base case (note that for $n = 0$ and $n = 1$ it holds trivially). Now, assume it holds for $n$. Show that it also holds for $n + 1$.
            
            We assume
            $$
                \langle n u, v \rangle = n \langle u, v \rangle
            $$

            Then,
            $$
                \langle (n + 1) u, v \rangle = \langle nu + u, v \rangle = \langle nu, v \rangle + \langle u, v \rangle = n \langle u, v \rangle + \langle u, v \rangle = (n + 1) \langle u, v \rangle
            $$

            This is what we wanted to show. Thus, by induction, we have shown homogeneity for all $\alpha \in \N$.

            Now, extend the property to $\Q$. Let $q \in \Q,\ q = \frac{n}{m}$.

            We can use the following trick:
            $$
                \langle u, v \rangle = \langle (\frac{m}{m} u), v \rangle = \langle m(\frac{1}{m} u), v \rangle = m \langle \frac{1}{m} u, v \rangle \overset{m \neq 0}\iff \frac{1}{m}\langle u, v \rangle = \langle \frac{1}{m} u, v \rangle
            $$

            Then, we have
            $$
                \langle \frac{n}{m} u, v \rangle = \langle n \cdot\frac{1}{m} u, v \rangle = n \langle \frac{1}{m} u, v \rangle = \frac{n}{m} \langle u, v \rangle
            $$

            So homogeneity holds for $\alpha \in \Q$ also. Finally, extend to $\R$. Because the function $f(\alpha) := \langle \alpha u, v \rangle$ is continuous (by continuity of the definition of the inner product here). Since $\Q$ is dense over $\R$, and homogeneity holds for all $q \in \Q$, it must hold for all $\alpha \in \R$ by the definition of continuity.

            So we have shown that homogeneity holds for all $\alpha \in \R$.

            We have additivity and homogeneity in the first variable, so this inner product is linear in the first variable.

            \item \textbf{Linearity in the second variable}: This can easily be shown using symmetry and linearity in the first variable.
        \end{itemize}

        We have shown that $\langle \cdot, \cdot \rangle$ as defined above is an inner product.

        Now, show that $\lVert v \rVert = \sqrt{\langle v, v \rangle}$.
        We have
        $$
            \begin{aligned}
                \langle v, v \rangle &= \frac{1}{2}\left(\lVert v + v \rVert^2 - \lVert v \rVert^2 - \lVert v \rVert^2\right) \\
                        &= \frac{1}{2}\left(\lVert 2v \rVert^2 - 2\lVert v \rVert^2\right) \overset{\text{homogeneity}}= \frac{1}{2}\left(4 \lVert v \rVert^2 - 2\lVert v \rVert^2\right) \\
                        &= \frac{1}{2}2 \lVert v \rVert^2 = \lVert v \rVert^2 \\
                        &\iff \sqrt{\langle v, v \rangle} = \lVert v \rVert
            \end{aligned}
        $$

        Thus, $\lVert v \rVert$ is induced by an inner product. \qed

    \end{itemize}

    \item Check that $\lVert \cdot \rVert_1$ defines a norm on $V$ and prove that it does not come from an inner product.
    
    Let us first show that $\lVert \cdot \rVert_1$ defines a norm on $V$.

    \begin{itemize}
        \item \textbf{Triangle inequality}: Fix $u, v \in V$. We want to show that
        $$
        \begin{aligned}
            \lVert u + v \rVert_1 \leq \lVert u \rVert_1 + \lVert v \rVert_1 &\iff \sum_{i = 1}^{n} |u_i + v_i| \leq \left( \sum_{i = 1}^{n} |u_i| \right) + \left( \sum_{i = 1}^{n} |v_i| \right)\\
        \end{aligned}
        $$

        By the standard triangle inequality on $\R$ (which we are allowed to assume), we have that
        $$
            \sum_{i = 1}^{n} |u_i + v_i| \leq \sum_{i = 1}^{n} \left(|u_i| + |v_i|\right) = \left( \sum_{i = 1}^{n} |u_i| \right) + \left( \sum_{i = 1}^{n} |v_i| \right)
        $$

        which is what we wanted to show. So the triangle inequality holds for $\lVert \cdot \rVert_1$.

        \item \textbf{Homogeneity}: Let $\alpha \in \R,\ v \in V$. We want to show that
        $$
            \lVert \alpha v \rVert_1 = |\alpha| \lVert v \rVert_1
        $$

        We have that
        $$
            \lVert \alpha v \rVert_1 = \sum_{i = 1}^{n} |\alpha v_i| = \sum_{i = 1}^{n} |\alpha| |v_i| = |\alpha| \sum_{i = 1}^{n} |v_i| = |\alpha| \lVert v \rVert_1
        $$

        which is what we wanted to show. So homogeneity holds for $\lVert \cdot \rVert_1$

        \item \textbf{Non-degeneracy}: Let $v \in V$ with $\lVert v \rVert_1 = 0$. Show that $v = 0$.
        We have that
        $$
            \lVert v \rVert_1 = 0 \iff \sum_{i = 1}^{n} |v_i| = 0
        $$

        Each $|v_i| \geq 0$, and the sum of non-negative numbers can only be zero if all are zero. So $v_i = 0\ \forall i \in \{1,\dots,n\} \iff v = 0$. So non-degeneracy holds.
    \end{itemize}

    Thus, we have shown that $\lVert \cdot \rVert_1$ defines a norm on $V$.

    Next, prove that $\lVert \cdot \rVert_1$ does not come from an inner product. We have already shown in (a) that a norm comes from an inner product iff it satisfies the parallelogram identity. Let us check if $\lVert \cdot \rVert_1$ satisfies the parallelogram identity. Try to find a counterexample in $V = \R^2$.

    Let $x = \begin{pmatrix} 1 \\ 0 \end{pmatrix},\ y = \begin{pmatrix} 0 \\ 1 \end{pmatrix}$. Then, the parallelogram identity tells us:
    $$
        \lVert x + y \rVert_1^2 + \lVert x - y \rVert_1^2 = 2 \lVert x \rVert_1^2 + 2 \lVert y \rVert_1^2
    $$
    $$
        \iff 2^2 + 2^2 = 2 \cdot 1 + 2 \cdot 1 \iff 8 = 4
    $$

    This does not hold, so $\lVert \cdot \rVert_1$ does not satisfy the parallelogram identity. Therefore, by the result from (a), $\lVert \cdot \rVert_1$ is not induced by an inner product. So $\lVert \cdot \rVert_1$ is a norm that is not induced by an inner product. \qed
\end{enumerate}

\begin{problem}{}{}
    Let $K = \R,\ V = M_{n \times n}(K)$, and consider the map
    $$
        V \times V \rightarrow K, \quad (A, B) \mapsto \mathrm{Tr}(A^TB)
    $$

    Show that it defines an inner product on $V$ and find an orthonormal basis with respect to this inner product. The induced norm is called the Hilbert-Schmidt norm. Give a formula for the norm of a matrix $A \in V$ in terms of its entries.
\end{problem}

First, show that this defines an inner product on $V$.

\begin{itemize}
    \item \textbf{Linearity in the first variable}: Let $A_1, A_2, B \in M_{n \times n}(\R)$. Using additivity of the trace, we can show that
    $$
        \mathrm{Tr}((A_1 + A_2)^T B) = \mathrm{Tr}((A_1^T + A_2^T) B) = \mathrm{Tr}(A_1^TB + A_2^TB) = \mathrm{Tr}(A_1^T B) + \mathrm{Tr}(A_2^T B)
    $$

    Next, show homogeneity. Fix $\alpha \in \R,\ A, B \in M_{n \times n}(\R)$. Using homogeneity of the trace, we have
    $$
        \mathrm{Tr}((\alpha A)^T B) = \mathrm{Tr}(\alpha A^T B) = \alpha \mathrm{Tr}(A^T B)
    $$

    So we have linearity in the first variable.

    \item \textbf{Symmetry}: Let $A, B \in M_{n \times n}(\R)$. Because the trace is just the sum of the diagonal, it is invariant under transposition. So
    $$
        \mathrm{Tr}(B^T A) = \mathrm{Tr}((B^T A)^T) = \mathrm{Tr}(A^T (B^T)^T) = \mathrm{Tr}(A^T B)
    $$

    So we have symmetry.

    \item \textbf{Linearity in the second variable}: Using symmetry and linearity in the first variable, we can easily show linearity in the second variable.
    \item \textbf{Positivity}: Let $A \in M_{n \times n}(\R)$ with $A \neq 0$. We want to show
    $$
        \mathrm{Tr}(A^T A) > 0
    $$

    Using the definitions, we get
    $$
        \mathrm{Tr}(A^T A) =: \sum_{i = 1}^{n} c_{ii}
    $$

    where 
    $$
        c_{ii} = \sum_{j = 1}^{n} A^T_{ij} \cdot a_{ji} = \sum_{j = 1}^{n} a_{ji} \cdot a_{ji} = \sum_{j = 1}^{n} a_{ji}^2
    $$

    where $a_{ji}$ is the entry of $A$ at the given index. Substituting back gives us
    $$
        \mathrm{Tr}(A^T A) = \sum_{i = 1}^{n} \sum_{j = 1}^{n} a_{ji}^2
    $$

    Since $A \neq 0$, at least one entry $a_{ji}$ must be non-zero. Then, $a_{ji}^2$ is $> 0$ (because our field is the real numbers). So it follows that
    $$
        \mathrm{Tr}(A^T A) = \sum_{i = 1}^{n} \sum_{j = 1}^{n} a_{ji}^2 > 0
    $$
    
    which is what we wanted to show.

\end{itemize}

So this does in fact define an inner product. Now, find an orthonormal basis with respect to this inner product. We are looking for $m := n \times n$ linearly independent matrices $(A_1,\dots,A_n)$ such that $\Span(A_1,\dots,A_m) = V = M_{n \times n}(\R)$ and $\lVert A_i \rVert = 1$ for all $1 \leq i \leq m$. We also require that $A_i \neq 0, A_j \neq 0$ and $\langle A_i, A_j \rangle = 0$ for all $1 \leq i \leq j \leq m,\ i \neq j$.

Let's see if we can just use the standard basis matrices $\mathcal{E} = (E_{11}, \dots, E_{nn})$ for $M_{n \times n}(\R)$. These are of course a basis for $M_{n \times n}(\R)$. Check if they are also orthonormal.

First, let's check normality. Let $E_{ij} \in \mathcal{E}$ be the matrix with zeros everywhere except at entry $(i, j)$. For this matrix, we have:
$$
    \lVert E_{ij} \rVert = \sqrt{\langle E_{ij}, E_{ij} \rangle} = \sqrt{\Tr \left(E_{ij}^T E_{ij}\right) } = \sqrt{\sum_{r = 1}^{n} \sum_{c = 1}^{n} (E_{ij})_{r, c}^2}
$$

A single entry of $E_{ij}$ is 1, all others are zero. So we have
$$
    \lVert E_{ij} \rVert = \sqrt{\sum_{r = 1}^{n} \sum_{c = 1}^{n} (E_{ij})_{r, c}^2} = 1
$$

Because $E_{ij}$ arbitrary, all $E_{ij} \in \mathcal{E}$ are normal. Show that they are also orthogonal. Let $E_{ij}, E_{kl} \in \mathcal{E}$ with $i \neq k$ or $j \neq l$. Show that they are orthogonal.

$$
    \langle E_{ij}, E_{kl} \rangle = \Tr(E_{ij}^T E_{kl}) = \sum_{r = 1}^{n} \sum_{c = 1}^{n} (E_{ij})_{r, c} \cdot (E_{kl})_{r, c}
$$

Because there is no index pair $(r, c)$ such that $(E_{ij})_{r, c} \neq 0$ and $(E_{kl})_{r, c} \neq 0$ simultaneously (because we fixed that they must be different elements of $\mathcal{E}$), we have that $(E_{ij})_{r, c} \cdot (E_{kl})_{r, c}$ is always zero. So
$$
    \langle E_{ij}, E_{kl} \rangle = \sum_{r = 1}^{n} \sum_{c = 1}^{n} (E_{ij})_{r, c} \cdot (E_{kl})_{r, c} = 0
$$

So we have shown that $\mathcal{E}$, the standard basis of $M_{n \times n}(\R)$ is an orthonormal basis with respect to the inner product we defined.

Finally, give a formula for the norm of a matrix $A \in V$ in terms of its entries. Fix $A \in V$. Then, as shown above,

$$
    \lVert A \rVert = \sqrt{\langle A, A \rangle} = \sqrt{ \Tr(A^T A) } = \sqrt{\sum_{i = 1}^{n} \sum_{j = 1}^{n} A_{ij}^2}
$$ \qed
\end{document}